\chapter{数值-形状双视角融合的多指标多节点异常检测}\label{chap:multi_metric}

本章针对分布式集群多指标多节点场景，研究数值与形状双视角融合的异常检测方法。首先介绍所使用的阶跃星辰GPU集群故障数据集及其特点；其次提出基于欧氏距离的数值特征异常检测算法；然后设计基于改进SAX的形状特征异常检测算法；最后通过加权融合两种视角的检测结果实现故障节点的精准定位，并在真实数据集上验证方法的有效性。

\section{集群多指标多节点数据集及异常检测算法框架}

本节首先对所使用的多指标多节点GPU集群异常检测数据集进行介绍，在此基础上分析多指标多节点异常检测问题的关键挑战，并提出数值与形状双视角融合的异常检测算法整体框架。

\subsection{多指标多节点数据集概述}

本节所使用的数据集为阶跃星辰GPU计算集群故障数据集，来源于真实生产环境中的大规模GPU计算集群监控系统，数据时间跨度从2024年7月至2025年1月。

大规模GPU集群作为支撑人工智能训练和高性能计算的核心基础设施，其稳定运行对业务连续性至关重要。集群通常由数百台GPU计算节点组成，通过统一的资源调度系统执行深度学习训练任务。当集群中某个节点发生故障时，不仅会导致训练任务失败，还可能因日志数据缺失而给后续根因分析带来更大挑战。尽管现有运维工具已能处理大部分常见故障，但对于少数难以定位的故障仍需投入大量人力进行排查，往往需要运维团队和算法团队各领域专家的经验才能解决问题。因此，基于多指标多节点的时间序列数据开展异常检测研究，对于实现GPU集群的智能化运维和故障快速定位具有重要的工程价值。

%% TODO: 插入图片 - 故障前后多节点时序对比图
%% 建议内容：展示故障发生前后各节点某一指标（如PCIe rx bytes或GPU利用率）的变化模式，
%% 突出显示故障节点与正常节点的行为差异，类似参考文档Figure 4的风格
% \begin{figure}[htbp]
%     \centering
%     \includegraphics[width=0.85\textwidth]{Img/multi_node_fault_comparison.pdf}
%     \caption{故障发生前后多节点时序对比示意图}
%     \label{fig:multi_node_fault_comparison}
% \end{figure}

该数据集包含105个故障案例作为测试集，同时构造了约100个无标签的正常运行作业作为训练集。每个故障案例的根因标注来自运维系统中的专家经验。表\ref{tab:fault_category}统计了测试集中各类故障的分布情况，涵盖GPU故障、机器故障、网络故障、PCIe故障及NCCL故障等多种类型。

\begin{table}[htbp]
    \centering
    \caption{故障类别及数量分布}
    \label{tab:fault_category}
    \begin{tabular}{lrl}
        \hline
        故障类别 & 数量 & 典型根因 \\
        \hline
        GPU故障 & 51 & Double bit error、XID error、GPU利用率不均衡 \\
        机器故障 & 13 & 节点宕机、机器检查错误 \\
        网络故障 & 18 & 交换机问题、光模块故障、线缆故障、信号质量差 \\
        PCIe故障 & 9 & PCIe性能退化、PCIe信号质量差 \\
        NCCL故障 & 5 & NCCL分段错误、NCCL超时 \\
        其他 & 9 & 未明确分类的故障 \\
        \hline
    \end{tabular}
\end{table}

在GPU集群运行过程中，监控系统通过周期性采样的方式持续采集各节点的运行状态数据，形成具有时间对齐关系的多节点多指标时间序列。每个节点会产生多个维度的监控指标，这些指标从不同角度反映节点的运行状态和硬件健康状况。根据指标数值是否随计算任务负载变化，可将监控指标分为负载相关指标和负载无关指标两类。负载无关指标（如PCIe Width、Network Link Down Event、Cable BER、XID Errors等）的数值不随计算任务变化，可直接通过与其他机器的对比找到离群节点；负载相关指标（如GPU利用率、GPU温度、显存使用率等）则需要更复杂的分析方法来区分正常波动与真实异常。表\ref{tab:monitoring_metrics}列出了数据集中的主要监控指标及其含义。

\begin{table}[htbp]
    \centering
    \caption{监控指标列表及说明}
    \label{tab:monitoring_metrics}
    \begin{tabular}{p{4.2cm}p{9.3cm}}
        \hline
        监控指标 & 说明 \\
        \hline
        CPU Usage & CPU繁忙处理任务的时间占比 \\
        Memory Usage & 系统RAM被程序和操作系统占用的比例 \\
        GPU PCIe Throughput & GPU与主板间通过PCIe传输数据的速率 \\
        NVLink Throughput & 通过高速NVLink互连传输数据的速率 \\
        GPU Temp & GPU的当前运行温度 \\
        Tensor Core Activity & GPU张量核心用于AI计算的利用率 \\
        FP32/FP16 Activity & 单精度/半精度浮点运算单元的活动百分比 \\
        SM Activity & GPU流式多处理器的利用率 \\
        GPU Memory Usage & GPU专用显存的使用比例 \\
        GPU Usage & GPU整体处理任务的繁忙时间占比 \\
        PCIe Width$^*$ & PCIe连接中活动数据通道的数量 \\
        Network Link Down Event$^*$ & 网络接口连接丢失的事件次数 \\
        Cable BER$^*$ & 线缆数据传输过程中的误码率 \\
        XID Errors$^*$ & GPU报告的严重错误次数 \\
        \hline
        \multicolumn{2}{l}{\footnotesize 注：带$^*$标记的为负载无关指标。}
    \end{tabular}
\end{table}

%% TODO: 插入图片 - 多指标时序数据示例图
%% 建议内容：展示同一时间段内多个不同指标的时序变化（如Mem Copy utilization、
%% Frame Buffer Used、SM clock、SM Occupancy、SM Activity等），
%% 体现数据的高维度、充满噪声等特点，类似参考文档第2页的多子图风格
% \begin{figure}[htbp]
%     \centering
%     \includegraphics[width=0.9\textwidth]{Img/multi_metric_timeseries.pdf}
%     \caption{多指标时序数据示例}
%     \label{fig:multi_metric_timeseries}
% \end{figure}

对于一个包含$N$个节点、$M$个监控指标、时间长度为$T$的数据集，可以形式化表示为一个三维张量：
\begin{equation}
    \mathbf{X} \in \mathbb{R}^{N \times M \times T}
\end{equation}

在实际运行环境中，集群节点通常执行相同类型的计算任务，并通过调度系统实现负载均衡，因此在正常运行情况下，各节点在相同指标上的变化趋势通常具有较高的一致性。当某个节点发生异常时，其监控指标曲线往往会偏离这种集群一致性模式。这种偏离可能表现为数值幅度上的异常，例如GPU温度显著高于其他节点；也可能表现为变化模式上的异常，例如功耗曲线出现异常波动或抖动，而数值幅度仍处于正常范围内。这种由系统调度机制和任务一致性所带来的节点间行为一致性，为基于相似性的异常检测方法提供了理论基础。

综上所述，该数据集具有以下特点：（1）数据来源于真实生产环境，具有较高的工程真实性和复杂性，不同节点之间存在一定程度的自然波动和噪声；（2）数据具有典型的多指标特征，不同指标对异常的敏感程度不同，仅依赖单一指标难以全面刻画节点异常行为；（3）数据具有显著的多节点关联特性，在正常运行状态下大多数节点表现出一致的运行模式，而异常节点则表现为相对于其他节点的偏离；（4）数据具有明显的时间序列特征，异常可能表现为短时突发异常，也可能表现为持续性异常。

\subsection{问题分析与算法框架}

在故障排查实践中存在一种朴素而有效的方法：找到与其他机器表现不同的机器，该机器大概率就是故障节点。基于这一思路，故障节点定位问题可以形式化定义为：给定一个包含$N$台机器的集合$M = \{m_1, m_2, \ldots, m_N\}$和一个不相似度函数$d(\cdot, \cdot)$，目标是找到机器$m^* \in M$，使得该机器与其他所有机器的累积不相似度最高，即：
\begin{equation}
    m^* = \arg \max_{m_k \in M} D(m_k)
\end{equation}
其中$D(m_k)$为机器$m_k$与集合中其他所有机器的总不相似度。

然而，基于多指标多节点时序数据的故障定位任务面临以下挑战：

\textbf{高度动态且变化的正常模式}。不同规模、不同类型任务的正常运行模式差异显著，传统异常检测方法往往只能识别时序数据中的主要模式，一旦存在其他正常模式变体就可能被误判为异常，导致较高的误报率。

\textbf{高维度且充满噪声的指标数据}。一方面节点的状态需要由多个不同来源的指标共同反映，另一方面这些指标自身因为传感器抖动或采样软件问题产生了大量噪声和抖动，通过这些数据准确刻画正常模式本身就十分困难。

\textbf{指标敏感度差异与融合策略}。不同指标对不同类型故障的敏感程度不同，需要设计合理的多指标融合策略，既能充分利用各指标的判别信息，又能避免单一指标的噪声干扰。

%% TODO: 插入图片 - 数值-形状双视角融合算法框架图
%% 建议内容：展示算法的整体流程框架，包含以下模块：
%% 1) 输入：多指标多节点时序数据
%% 2) 数值视角分支：欧氏距离计算 -> 多指标加权融合 -> 数值异常分数
%% 3) 形状视角分支：改进SAX符号化 -> 符号序列距离 -> 形状异常分数  
%% 4) 融合模块：加权融合 -> 最终异常分数 -> 故障节点排序
% \begin{figure}[htbp]
%     \centering
%     \includegraphics[width=0.95\textwidth]{Img/dual_perspective_framework.pdf}
%     \caption{数值-形状双视角融合异常检测框架}
%     \label{fig:dual_perspective_framework}
% \end{figure}

针对上述挑战，本章提出数值与形状双视角融合的异常检测框架，其核心思想是从数值特征和形状特征两个互补视角分析节点异常行为，然后通过加权融合得到最终的异常判定结果。该框架主要包含三个阶段：

\textbf{数值特征异常检测阶段}。采用基于欧氏距离的相似度度量方法，通过计算节点间时间序列的数值距离来量化差异程度。对于每个监控指标，计算各节点与其他节点的平均距离作为该指标下的异常分数，然后采用加权平均策略融合多个指标的异常分数，得到数值视角下的综合异常评分。

\textbf{形状特征异常检测阶段}。采用改进的SAX符号化算法，通过将时间序列映射为符号序列并计算符号序列间的距离来捕捉曲线形态的差异。该方法取消了传统SAX的分段聚合步骤以保留原始时间分辨率，并采用严格距离计算规则以提升对形状变化的敏感性。

\textbf{多视角融合与故障定位阶段}。将两个视角的异常分数进行加权融合，综合考虑数值偏离和形态变化两方面因素，输出节点的最终异常分数并进行排序。排名靠前的节点即为最可能发生故障的候选节点，供运维人员进一步排查确认。

\section{基于数值特征的异常检测算法}

在分布式集群的性能监控中，节点的多个监控指标在正常状态下往往呈现出与集群整体一致的行为模式。当某个节点发生故障时，其一个或多个监控指标的数值会偏离集群的正常范围，表现为与其他节点的显著差异。基于这一观察，本节提出一种基于数值特征的异常检测算法，其核心思想是：对每个监控指标独立计算节点间的欧氏距离以量化数值差异，通过Z-score标准化将距离转换为可比较的异常分数，最后采用加权融合策略整合多指标信息，输出节点的综合异常评分。

%% TODO: 插入图片 - 数值特征异常检测算法流程图
%% 建议内容：展示算法的核心流程，包含以下步骤：
%% 1) 输入：多节点多指标时序数据 X ∈ R^{N×M×T}
%% 2) 单指标处理：对每个指标m，计算N×N距离矩阵 → 行求和得到距离向量 → Z-score标准化得到异常分数向量
%% 3) 多指标融合：加权求和各指标异常分数 → 输出最终异常分数
%% 可采用流程图或数据流图的形式，突出"先分后合"的处理策略
% \begin{figure}[htbp]
%     \centering
%     \includegraphics[width=0.9\textwidth]{Img/numerical_detection_pipeline.pdf}
%     \caption{基于数值特征的异常检测算法流程}
%     \label{fig:numerical_detection_pipeline}
% \end{figure}

\subsection{基于欧氏距离的单指标异常度量}

在多指标多节点场景中，一个直观的思路是将所有指标的时序数据拼接成高维向量，直接计算节点间的多指标联合距离。然而，这种联合计算方式存在以下问题：

\textbf{量纲与尺度差异}。不同监控指标的数值范围和单位存在显著差异，例如CPU利用率的取值范围为0\%-100\%，而PCIe吞吐量可能达到数十GB/s。直接进行多维欧氏距离计算时，数值尺度较大的指标会主导距离计算结果，导致其他指标的异常信息被掩盖。

\textbf{故障敏感度差异}。不同类型的故障往往只影响部分相关指标，而非所有指标同时异常。例如，GPU温度异常可能仅导致温度和功耗指标偏离正常范围，而内存使用率和网络吞吐量可能保持正常。若采用多指标联合距离，正常指标的距离贡献会稀释异常指标的判别信号，降低检测灵敏度。

\textbf{可解释性需求}。在实际运维场景中，不仅需要识别异常节点，还需要定位导致异常的具体指标，以指导后续的故障排查。联合距离计算将多指标信息融合为单一数值，丧失了各指标的独立贡献信息。

基于上述考虑，本文采用"先分后合"的策略：首先对每个监控指标独立计算节点间的欧氏距离并转换为异常分数，然后通过加权融合整合多指标的异常信息。这种策略的优势在于：（1）各指标在独立空间中进行距离计算和标准化，避免了量纲差异的干扰；（2）每个指标的异常分数独立可查，便于定位故障相关指标；（3）通过权重调整可以灵活控制各指标对最终结果的贡献程度。

对于单个监控指标$m$，设集群中$N$个节点在时间窗口$[1,T]$内的时序数据为$\mathbf{X}^{(m)} \in \mathbb{R}^{N \times T}$，其中$\mathbf{x}_i^{(m)} = (x_{i,1}^{(m)}, x_{i,2}^{(m)}, \ldots, x_{i,T}^{(m)})$表示节点$i$在指标$m$上的时间序列。节点$i$与节点$j$之间的欧氏距离定义为：
\begin{equation}
    d_{ij}^{(m)} = \sqrt{\sum_{t=1}^{T} (x_{i,t}^{(m)} - x_{j,t}^{(m)})^2}
\end{equation}

通过计算所有节点对之间的距离，构建距离矩阵$\mathbf{D}^{(m)} \in \mathbb{R}^{N \times N}$。节点$i$在指标$m$上的累积距离定义为该节点与其他所有节点距离的平均值：
\begin{equation}
    \bar{d}_i^{(m)} = \frac{1}{N-1} \sum_{j \neq i} d_{ij}^{(m)}
\end{equation}

累积距离越大，表明该节点在该指标上与集群整体的偏离程度越高。为了将不同指标的累积距离转换为可比较的异常分数，采用Z-score标准化方法：
\begin{equation}
    s_i^{(m)} = \frac{\bar{d}_i^{(m)} - \mu_d^{(m)}}{\sigma_d^{(m)}}
\end{equation}
其中$\mu_d^{(m)}$和$\sigma_d^{(m)}$分别为所有节点累积距离的均值和标准差。标准化后的异常分数$s_i^{(m)}$表示节点$i$在指标$m$上的累积距离相对于集群平均水平的偏离程度，正值越大表示异常程度越高。

\subsection{多指标加权融合策略}

在获得各指标的独立异常分数后，需要将其融合为节点的综合异常评分。本文采用加权求和的融合策略，节点$i$的综合异常分数定义为：
\begin{equation}
    S_i = \sum_{m=1}^{M} w_m \cdot \max(s_i^{(m)}, 0)
    \label{eq:weighted_fusion}
\end{equation}
其中$w_m$为指标$m$的权重，满足$\sum_{m=1}^{M} w_m = 1$。

公式\ref{eq:weighted_fusion}中对异常分数取$\max(\cdot, 0)$操作的原因如下：当$s_i^{(m)} < 0$时，表示节点$i$在指标$m$上的累积距离低于集群平均水平，即该节点在该指标上表现正常甚至优于平均。在故障检测场景中，这种"正常"信号不应对最终的异常判定产生负面影响，因此将负值截断为0。此外，当某指标的异常分数为0时，意味着该指标在当前时间段未能提供有效的异常判别信息，在融合时自然不产生贡献。

权重的设定可以采用以下策略：

\textbf{等权重策略}。在缺乏先验知识的情况下，对所有指标赋予相等权重，即$w_m = 1/M$。这种策略简单直观，适用于探索性分析阶段。

\textbf{基于领域知识的权重}。根据运维经验，对关键性指标（如GPU利用率、温度、功耗等）赋予较高权重，对辅助性指标赋予较低权重。这种策略能够突出重要指标的判别作用。

\textbf{基于历史数据的自适应权重}。通过分析历史故障案例中各指标的检测效果，自动学习最优权重配置。例如，可以采用基于验证集性能的网格搜索或贝叶斯优化方法确定权重。

在本文的实验中，默认采用等权重策略，并在消融实验中验证不同权重设置对检测性能的影响。

\subsection{算法实现与效率优化}

算法\ref{alg:numerical_detection}给出了基于数值特征的异常检测算法的完整流程。算法输入为多节点多指标时序数据张量$\mathbf{X} \in \mathbb{R}^{N \times M \times T}$和指标权重向量$\mathbf{w}$，输出为各节点的综合异常分数$\mathbf{S}$。

\begin{algorithm}[htbp]
    \caption{基于数值特征的多指标异常检测算法}
    \label{alg:numerical_detection}
    \begin{algorithmic}[1]
        \Require 时序数据 $\mathbf{X} \in \mathbb{R}^{N \times M \times T}$，权重向量 $\mathbf{w} \in \mathbb{R}^M$
        \Ensure 异常分数向量 $\mathbf{S} \in \mathbb{R}^N$
        \State 初始化 $\mathbf{S} \gets \mathbf{0}$
        \For{$m = 1$ to $M$} \Comment{遍历每个指标}
            \State $\mathbf{D}^{(m)} \gets$ 计算节点间欧氏距离矩阵 \Comment{$\mathbf{D}^{(m)} \in \mathbb{R}^{N \times N}$}
            \State $\bar{\mathbf{d}}^{(m)} \gets$ 计算各节点累积距离 \Comment{$\bar{d}_i^{(m)} = \frac{1}{N-1}\sum_{j \neq i} D_{ij}^{(m)}$}
            \State $\mathbf{s}^{(m)} \gets$ Z-score标准化$(\bar{\mathbf{d}}^{(m)})$ \Comment{转换为异常分数}
            \State $\mathbf{S} \gets \mathbf{S} + w_m \cdot \max(\mathbf{s}^{(m)}, 0)$ \Comment{加权累加}
        \EndFor
        \State \Return $\mathbf{S}$
    \end{algorithmic}
\end{algorithm}

该算法的时间复杂度为$O(M \cdot N^2 \cdot T)$，其中距离矩阵计算是主要的计算开销。为满足大规模集群实时监控的需求，本文在算法实现层面进行了以下效率优化：

\textbf{向量化距离计算}。采用NumPy和SciPy库的向量化计算功能，将循环操作转换为矩阵运算。具体地，利用scipy.spatial.distance.cdist函数一次性计算整个距离矩阵，避免显式的双重循环，充分利用底层优化的线性代数库（如BLAS）的并行计算能力。

\textbf{数据预处理}。对监控数据进行预处理以确保数据质量：采用线性插值填充缺失值，确保时序数据的完整性；按时间窗口切分数据，将长时序分解为可独立处理的短窗口，降低单次计算的内存占用。

\textbf{增量计算与缓存}。对于滑动窗口检测场景，采用增量更新策略：当新数据到达时，仅重新计算涉及新数据点的距离，而非完全重建距离矩阵。同时缓存中间计算结果（如各节点的历史累积距离统计量），减少重复计算。

通过上述优化措施，算法在处理包含50个节点、6个监控指标、1000个时间点的典型数据规模时，单次检测耗时可控制在0.3秒以内，满足秒级响应的实时监控需求。

\section{基于形状特征的异常检测算法}

\subsection{经典SAX算法原理}

SAX是一种经典的时间序列符号化表示方法，其核心思想是将连续的时间序列转换为离散的符号序列，从而实现时间序列的降维和模式匹配。传统SAX算法主要包括四个步骤。

第一步是归一化处理。将原始时间序列标准化为均值为0、标准差为1的标准正态分布，以消除量纲和基线差异。对于时间序列$\{x_1, x_2, \ldots, x_T\}$，归一化后的序列为：
\begin{equation}
    \hat{x}_t = \frac{x_t - \mu}{\sigma}
\end{equation}
其中$\mu$和$\sigma$分别为序列的均值和标准差。

第二步是分段聚合近似。将归一化后的时间序列划分为$w$个等长的片段，并计算每个片段内所有数据点的平均值作为该片段的代表值。这一步骤实现了序列的降维，将长度为$T$的序列压缩为长度为$w$的表示。

第三步是符号化映射。基于标准正态分布的分位数点将数值空间划分为$\alpha$个等概率的区间，其中$\alpha$为字母表的大小。每个片段的平均值根据其所落入的区间被映射到对应的符号。通过这种方式，连续的时间序列被转换为由$\alpha$个符号组成的符号序列。

第四步是距离计算。在符号空间中定义符号间的距离，进而计算符号序列之间的距离。传统SAX距离的定义中，相邻符号之间的距离为0，非相邻符号之间的距离为对应分位数点之差。

\subsection{基于SAX算法的符号化表示与模式匹配}

为适应GPU集群故障检测场景的特殊需求，本文对传统SAX算法进行了两方面改进。

第一个改进是取消分段聚合步骤，保留原始时间分辨率。传统SAX的分段聚合虽然实现了降维，但可能丢失重要的细节信息。在GPU性能监控场景中，瞬时尖峰和短暂抖动往往包含重要的故障指示信息。因此，本文取消了分段处理，对归一化后的每个数据点直接进行符号化，使算法能够捕捉到更细微的时间点异常。

第二个改进是采用严格距离计算规则。传统SAX算法中，相邻符号被视为相似，其距离设为0。在故障检测场景中，这种宽容性可能导致对形状变化的不敏感从而造成漏报。本文修改距离计算规则，仅当两个符号完全相同时距离才为0，否则计算实际距离：
\begin{equation}
    d'(l_i, l_j) = \begin{cases}
        0, & \text{if } i = j \\
        \beta_{\max(i,j)-1} - \beta_{\min(i,j)}, & \text{if } i \neq j
    \end{cases}
\end{equation}
其中$\beta_k$为第$k$个高斯分位数点。

基于改进的SAX算法，多节点形状异常检测的流程如下。首先对各节点的时间序列进行归一化处理，然后根据预设字母表大小计算高斯分位数点并将归一化序列映射为符号序列。接着计算所有节点对之间的符号序列距离，构建节点距离矩阵。最后计算每个节点与其他节点的平均距离，基于Z-score方法计算异常分数，并对各指标的异常分数进行加权聚合得到最终结果。

\subsection{算法效率优化}

为满足大规模GPU集群实时监控的需求，本文在算法实现层面进行了效率优化。

在距离矩阵计算方面，采用预计算符号距离表和向量化操作来避免运行时的重复计算。对于字母表大小为$\alpha$的SAX表示，预先计算并存储$\alpha \times \alpha$的符号距离矩阵，在运行时通过查表方式获取距离值。同时利用NumPy和SciPy库的向量化计算功能，将循环操作转换为矩阵运算，充分利用底层优化的线性代数库。

在数据处理方面，采用批量处理策略和内存优化技术。对于大规模数据集，按时间窗口进行分批处理，避免一次性加载全部数据导致的内存压力。同时采用稀疏矩阵表示和即时计算策略，仅在需要时计算特定节点对之间的距离，减少不必要的计算开销。

通过上述优化措施，算法在处理包含数百个节点、数千个时间点的多指标数据集时，单次检测耗时能够控制在秒级，满足实时监控的响应要求。

\section{多算法融合与故障定位}

\subsection{加权融合设计}

在多算法融合阶段，需要将基于数值特征和基于形状特征的检测结果进行综合。本文采用加权平均的融合策略，具体设计如下。

设节点$i$通过数值检测算法得到的异常分数为$s_i^{num}$，通过形状检测算法得到的异常分数为$s_i^{shape}$。由于两种算法输出的分数尺度可能不同，首先对各自的分数进行归一化处理，使其落入可比较的范围。然后通过加权平均得到融合后的异常分数：
\begin{equation}
    s_i^{final} = w_{num} \cdot s_i^{num} + w_{shape} \cdot s_i^{shape}
\end{equation}
其中$w_{num}$和$w_{shape}$分别为数值检测和形状检测的权重，满足$w_{num} + w_{shape} = 1$。

权重的设置可以基于领域知识或历史检测性能进行调整。在默认情况下，本文采用等权重设置，即$w_{num} = w_{shape} = 0.5$。在实际应用中，如果已知某类故障更多表现为数值异常，可以适当提高数值检测的权重，反之亦然。

对于多指标场景，在单一视角内部也需要进行多指标融合。本文采用的策略是对各指标分别计算异常分数，然后进行加权聚合。在加权聚合时，遵循以下规则以提高鲁棒性。异常分数为0的指标不参与聚合计算，因为这表明该指标在当前时间段内未能提供有效的异常信息。异常分数为负数的指标按0计算，因为负的Z-score意味着该节点在该指标上表现正常甚至优于平均水平，不应对最终的异常判定产生负面影响。

\subsection{故障定位与根因分析建议}

在获得各节点的融合异常分数后，根据分数对节点进行降序排序，排名靠前的节点即为最可能发生故障的节点。在实际应用中，通常取Top-K个节点作为故障候选进行进一步排查。

基于数值和形状双视角的检测结果，还可以为故障类型判断提供参考。当节点表现出高数值异常分数且高形状异常分数时，可能对应严重的硬件故障或系统级问题，例如GPU核心损坏或散热系统失效。当节点表现出高数值异常分数但低形状异常分数时，可能对应数值漂移或配置异常，例如功耗阈值设置不当或资源分配失衡。当节点表现出低数值异常分数但高形状异常分数时，可能对应间歇性故障或软件异常，例如驱动程序问题或进程异常。

这种故障类型的初步判断能够为运维人员的后续排查工作提供方向性指导，提高故障定位的效率。

\section{实验验证与对比分析}

\subsection{实验设置}

本节实验在阶跃星辰GPU计算集群故障数据集上进行。该数据集包含105个真实故障案例，每个案例对应一次明确记录的GPU节点故障事件，故障类型涵盖GPU性能退化、温度异常升高、功耗异常波动等。在每个故障案例中，数据集包含故障节点及其所在集群所有节点在故障发生时间窗口内的多指标时间序列数据。

实验评估采用Acc@5和Avg Rank两个指标。Acc@5表示真实故障机器出现在算法检测Top-5中的案例比例，衡量算法的故障检测准确率。Avg Rank为真实故障机器在Top-5中的平均排名，衡量算法的故障定位效率，数值越小表示定位越精准。

基线算法选择马氏距离方法，该方法考虑了变量之间的协方差关系，是多元异常检测领域的经典方法。

\subsection{方法结果对比分析与案例研究}

表\ref{tab:multi_metric_results}展示了各方法在105个故障案例上的检测效果对比。

\begin{table}[htbp]
    \centering
    \caption{多指标多节点异常检测方法对比}
    \label{tab:multi_metric_results}
    \begin{tabular}{lcc}
        \hline
        算法 & Acc@5 & Avg Rank \\
        \hline
        马氏距离 & 0.453 & 1.20 \\
        欧氏距离+加权聚合 & 0.800 & 1.25 \\
        改进SAX+加权聚合 & 0.762 & 1.30 \\
        数值-形状融合 & 0.838 & 1.18 \\
        \hline
    \end{tabular}
\end{table}

从结果可以看出，基于欧氏距离的数值检测方法在Acc@5指标上达到0.800，显著优于马氏距离基线方法的0.453。这表明在GPU集群场景中，基于节点间曲线相似性的检测思路比基于变量协方差关系的传统方法更为有效。

基于改进SAX的形状检测方法取得了0.762的Acc@5，虽然略低于数值检测方法，但能够捕捉到一些数值方法遗漏的形态异常案例。将两种方法进行融合后，Acc@5提升至0.838，Avg Rank降至1.18，验证了双视角融合策略的有效性。

\subsection{融合方法优势分析}

融合方法相比单一视角方法具有两方面优势。

第一是更全面的异常覆盖。数值检测方法擅长识别幅度显著偏离的异常，而形状检测方法擅长识别变化模式异常但幅度正常的情况。通过融合两种视角，能够覆盖更多类型的异常模式。实验统计表明，约15\%的故障案例主要表现为形态异常而非数值异常，单纯依赖数值方法会造成漏报。

第二是更稳定的检测性能。单一方法可能在特定类型的故障上表现优异但在其他类型上效果欠佳，而融合方法通过综合多个视角的判断结果，能够获得更稳定的整体性能，降低了对特定故障类型的依赖。

\subsection{算法效率实验对比}

为验证算法效率优化的效果，本节对比了优化前后的算法执行时间。实验选取包含50个节点、6个监控指标、1000个时间点的典型数据规模进行测试。

\begin{table}[htbp]
    \centering
    \caption{算法效率对比}
    \label{tab:efficiency}
    \begin{tabular}{lcc}
        \hline
        算法 & 平均执行时间 & 加速比 \\
        \hline
        欧氏距离（优化前） & 2.35s & 1.0$\times$ \\
        欧氏距离（优化后） & 0.28s & 8.4$\times$ \\
        改进SAX（优化前） & 5.12s & 1.0$\times$ \\
        改进SAX（优化后） & 0.45s & 11.4$\times$ \\
        \hline
    \end{tabular}
\end{table}

通过向量化计算和预计算优化，欧氏距离方法的执行时间从2.35秒降至0.28秒，加速比达到8.4倍。改进SAX方法的执行时间从5.12秒降至0.45秒，加速比达到11.4倍。优化后的算法能够满足实时监控场景的响应要求。

\section{本章小结}

本章针对多指标多节点场景下的GPU集群异常检测问题，提出了数值与形状双视角融合的异常检测框架。在数值视角方面，采用基于欧氏距离的多指标加权融合方法量化节点间的数值差异。在形状视角方面，提出了改进的SAX符号化表示算法，通过保留原始时间分辨率和采用严格距离计算规则来提升对形状变化的敏感性。最后通过加权融合两种视角的检测结果，实现了对不同类型异常的全面覆盖。在真实GPU集群故障数据集上的实验表明，所提融合方法的故障检测准确率达到83.8\%，显著优于单一视角方法和传统基线方法。
